First, we focus on simple rulesets for low PD-L1 expression thresholds.
High values of $Q_y$ led to very little learning as discussed in \autoref{sec:boolean-results}.
Therefore, we restrict ourselves to $Q_y = 0.5, 0.6.$
We are interested in the performance of very simple rulesets and what they look like.
Therefore, we study
\[C = 1,2 \quad \text{and} \quad D = 1,2,3,\]
training on Broad in full.
We fix $\nq = 9$ and consider both lower-bounds and upper-bounds.
The use of the exact pricer has proven to be time-consuming, especially for higher $\nq$.
We are interested in high $\nq$ to rule out the effect of course discretization on ruleset performance.
To avoid wasting time and compute (that could otherwise be used by other users of the HPC),
we use only the heuristic pricer with beam widths $(1000,1000)$ for $D=2$ and $(1000,1000,1000)$ for $D=3.$
While this removes optimality guarantees like those discussed in \autoref{sec:boolean_model_appendix},
rulesets still provide a proof that a certain performance can be attained.
Furthermore, optimal solutions for small $C$ and small $D$ will be discussed in \autoref{sec:threshold_greedy}.

The total time limit is set to 3 hours, which is much more than needed for this setup.

Next, we study the effect of increasing $C$ and $D$ on performance.
Again, Broad is used in full for training.
The settings for this experiment are analogous to the experiment above.

Finally, we study the generalization of simple rulesets for $Q_y=0.5, 0.6.$
We apply 10-fold stratified cross-validation on Broad.
Then we compare train and test accuracy.
We set $D = 2$ and study $C=1,2,3$.
We also study the consistency of these solutions by counting how often each ruleset occurs among splits.
In this count, we disregard the threshold value and direction (lower-bound or upper-bound) of each clause,
but identify a ruleset only by which TFs appear together in each rule.
Ideally, the ruleset found is consistent, as it is meant to be a general rule governing PD-L1 expression.